\section{Extracting a choice from a bounded-bag witness}
\label{sec:separator}

The dual equations force odd total weight at every type.
We now use small separators to locate one bag whose main vertices
give compatible choices.
The odd-side orientation below is closely related to the separator
argument for oddomorphisms of Neuen and Seppelt
\citep[Lemma~4.4 and proof of Theorem~4.1]{NS2026}.
Here the helper types additionally force compatibility of the
representatives; we retain the local orientation to control the
weighted contribution of every outer side.

\begin{lemma}[Choice extraction]\label{lem:extract-choice}
Let \(F\) have a tree decomposition with bags of size at most \(t\).
Suppose that a homomorphism \(\phi:F\to G\) and weights
\(s:V(F)\to\FF\) and \(r:E(F)\to\FF\) satisfy
\eqref{eq:dual}, where \(a=\tau\circ\phi\).
Then the choice system has a successful choice.
\end{lemma}

\begin{proof}
All sums in this proof are in \(\FF\).
We use the type masses \(m_u(Y)\) from \eqref{eq:type-mass};
in particular, \(m_u(V(F))=1\) for every type \(u\).

\smallskip
\noindent\emph{Orienting the decomposition.}
Fix a decomposition tree \(\mathcal D\), and contract edges joining
identical bags.  Every remaining edge has an intersection \(S\) of
size at most \(t-1\): an intersection of size \(t\) would force both
incident bags to be that same set.
Deleting the tree edge partitions \(V(F)\) into \(S\) and two sides
\(A^0,A^1\).  Indeed, the occurrence subtree of a vertex outside \(S\)
lies in just one of the two tree components.  No graph edge joins the
two sides, since its endpoints occur together in a bag.

At least one main type is absent from \(S\).
If two main types \(u,v\) are absent, all \(uv\)-edges incident with
either type on a side stay on that side.  Summing \eqref{eq:dual}
over their two endpoint classes gives
\(m_u(A^\nu)=m_v(A^\nu)\) for \(\nu\in\{0,1\}\).
For any absent main type \(u\),
\[
m_u(A^0)+m_u(A^1)=m_u(V(F))=1.
\]
Orient the tree edge toward the side of mass \(1\).
This is well defined whether there is one absent main type or several.

Choose a sink of the oriented tree and denote its bag by \(B\).
For each incident edge let \(A\) be its outer side and \(S\subseteq B\)
its intersection.  These outer sides partition \(V(F)\setminus B\).
To see this precisely, the occurrence subtree of a vertex outside
\(B\) avoids the sink node and lies in exactly one component of
\(\mathcal D\) minus that node.  A vertex in \(B\) occurring in a branch
must belong to that branch's intersection \(S\), so is excluded from
its outer side.  Thus vertices are counted once, regardless of how
many bags contain them.  Moreover, every neighbor of \(A\) outside
\(A\) lies in \(S\), and the sink orientation gives
\[
m_\ell(A)=0
\qquad\text{for every main type \(\ell\) absent from \(S\).}
\]
If a main type were absent from \(B\), it would be absent from every
such \(S\) and have mass zero in the bag and in all outer sides,
contrary to its global mass \(1\).
Since \(|B|\le t\), the bag therefore consists of exactly one vertex
\(x_i\) of each main type \(i\), and no helper vertices.
No condition on the individual weights \(s(x_i)\) is needed.

\smallskip
\noindent\emph{A global mismatch test.}
Write \(\phi(x_i)=(i,d_i)\), with \(d_i\in D_i\).
The bag \(B\) need not induce a clique, so compatibility of these
representatives does not follow from \(\phi\) alone.
Suppose, for a contradiction, that some pair \(i<j\) has different
labels, and put \(\alpha=\lambda_{ij}^i(d_i)\).
We choose label predicates that vanish on these representatives and are
complementary on matching labels. For source vertices of types \(i,j\),
respectively, put
\[
p_i(x)=\ind{\lambda_{ij}^i(\phi(x))\ne\alpha},
\qquad
p_j(x)=\ind{\lambda_{ij}^j(\phi(x))=\alpha}.
\]
For \(Y\subseteq V(F)\), let
\[
P(Y)=\sum_{d\in\{i,j\}}\ \sum_{\substack{x\in Y\\a(x)=d}}
                   s(x)p_d(x).
\]
Each \(ij\)-edge has matching labels under \(\phi\), hence the
two endpoint predicates sum to \(1\).  Expanding the corresponding
coordinates of \eqref{eq:dual} gives
\begin{equation}\label{eq:global-mismatch}
\begin{aligned}
P(V(F))
 &=\sum_{\substack{xy\in E(F)\\a(x)=i,\ a(y)=j}}
          r(xy)\bigl(p_i(x)+p_j(y)\bigr)\\
 &=\sum_{\substack{xy\in E(F)\\a(x)=i,\ a(y)=j}}r(xy)
   =m_i(V(F))=1.
\end{aligned}
\end{equation}
Here each edge is listed with its \(i\)-endpoint first.
On the other hand, \(P(B)=0\).

\smallskip
\noindent\emph{Vanishing on an outer side.}
Fix an outer side \(A\) with intersection \(S\), and choose a main type
\(\ell\) absent from \(S\); thus \(m_\ell(A)=0\).
We show that \(P(A)=0\), retaining the possible boundary edges.

First suppose \(\ell=i\); the case \(\ell=j\) is symmetric.
Every \(i\)-neighbor of a \(j\)-vertex in \(A\) lies in \(A\),
because an external neighbor would lie in \(S\), which has no
\(i\)-vertex.  An edge from an \(i\)-vertex in \(A\) may reach a
\(j\)-vertex in \(S\), but its \(p_j\)-value is zero.
Consequently expansion of \eqref{eq:dual} yields
\[
\begin{aligned}
P(A)
 &=\sum_{\substack{x\in A\\a(x)=i}}
       \ \sum_{\substack{y\in N_F(x)\\a(y)=j}}
         r(xy)\bigl(p_i(x)+p_j(y)\bigr)\\
 &=\sum_{\substack{x\in A\\a(x)=i}}
       \ \sum_{\substack{y\in N_F(x)\\a(y)=j}}r(xy)
   =m_i(A)=0.
\end{aligned}
\]

Now suppose \(\ell\notin\{i,j\}\), and put \(h=h_{ij;\ell}\).
Each base vertex of type \(h\) couples the two predicates through its
stored \(L_{ij}\)-label. Having this helper type for every possible \(\ell\)
makes the calculation available whichever main type \(S\) omits.
Since \(S\subseteq B\), neither \(h\) nor \(\ell\) occurs in \(S\).
We will show \(P(A)=m_h(A)=m_\ell(A)=0\), using a zero boundary
correction to complete the neighbor sums in the first equality.
Write \(H_A=\{w\in A:a(w)=h\}\).
Every \(h\)-neighbor of an \(i\)- or \(j\)-vertex in \(A\) belongs to
\(H_A\): an outside neighbor would have to lie in \(S\).
Expanding the \(h\)-coordinate at each tested main vertex therefore gives
\[
P(A)=\sum_{w\in H_A}\ \sum_{d\in\{i,j\}}
       \ \sum_{\substack{x\in A\cap N_F(w)\\a(x)=d}}
                  r(wx)p_d(x).
\]
The boundary correction
\[
C_S=\sum_{w\in H_A}\ \sum_{d\in\{i,j\}}
       \ \sum_{\substack{x\in S\cap N_F(w)\\a(x)=d}}
                  r(wx)p_d(x)
\]
is zero, since \(p_i,p_j\) vanish on their respective vertices of \(S\).
Adding it includes \emph{all} \(i\)- and \(j\)-neighbors of every
\(w\in H_A\), because all external neighbors of \(A\) lie in \(S\).

For each such \(w\), let \(q_w\in L_{ij}\) be the label stored at
\(\phi(w)\), and write
\[
R_d(w)=\sum_{\substack{x\in N_F(w)\\a(x)=d}}r(wx)
       =s(w)\qquad(d=i,j),
\]
where the equality is the corresponding coordinate of
\eqref{eq:dual}.
Its \(i\)-neighbors have predicate \(\ind{q_w\ne\alpha}\),
and its \(j\)-neighbors have predicate \(\ind{q_w=\alpha}\).
Hence
\[
\begin{aligned}
P(A)
 &=P(A)+C_S\\
 &=\sum_{w\in H_A}
       \bigl(\ind{q_w\ne\alpha}R_i(w)
               +\ind{q_w=\alpha}R_j(w)\bigr)\\
 &=\sum_{w\in H_A}s(w)=m_h(A).
\end{aligned}
\]
These are separate coordinate equations at each source vertex \(w\);
its number of neighbors, its label \(q_w\), and its weight \(s(w)\)
may vary.  An empty neighbor class simply forces \(s(w)=0\).
Finally, because neither \(h\) nor \(\ell\) occurs in \(S\), every
\(h\ell\)-edge incident to either type in \(A\) stays in \(A\).
Summing its coordinate equations at the two endpoint classes gives
\(m_h(A)=m_\ell(A)=0\), as required.

The disjoint vertex partition now gives
\[
P(V(F))=P(B)+\sum_{\text{outer sides }A}P(A)=0,
\]
contradicting \eqref{eq:global-mismatch}; the sum is empty if the
decomposition has only one bag.
Thus all \(d_i\) have matching pair labels and form a successful choice.
\end{proof}

\begin{proof}[Proof of Theorem~\ref{thm:choice-reduction}]
The parity lifts in Definition~\ref{def:lifts} are simple, uncolored and
of equal order; their preceding tag count gives \(n\le A M_t\).
If a successful choice exists, Lemma~\ref{lem:choice-witness} gives a
strict count difference for \(T_t\). Lemma~\ref{lem:test-detected}
detects that difference after two joint rounds or one separate round.

Conversely, suppose that no successful choice exists. If a graph \(F\)
with bags of size at most \(t\) had different counts into the two lifts,
Lemma~\ref{lem:lift-duality} would give the data \(\phi,s,r\) of
Lemma~\ref{lem:extract-choice}. That lemma would produce a successful
choice, a contradiction. All such homomorphism counts therefore agree.
Lemma~\ref{lem:hom-implies-wl} gives equivalence in the corresponding
convention, including the pure-tuple disjoint-union implementation.
\end{proof}
